\chapter{Introduction}\label{ch:introduction}

The Bak-Tang-Wiesenfeld (BTW) sandpile model is a classical example that demonstrates self-organised criticality (SOC), a property of certain systems evolving towards a statistically stationary critical state without finetuning the initial conditions \cite{BTW1988}.  \\

%Q1- why scale invariance described by power laws
A system that exhibits SOC are expected to show scale invariant properties. Being scale invariant means that,
\[
p(cx)=g(c)p(x)
\]
where $c$ is the scaling factor. A power law satisfies exactly this property as:
\[
p(x)= Ax^{-\tau}
\]
and hence,
\[
p(cx) = A(cx)^{-\tau}=g(c)p(x)
\]
where $g(c)=c^{-\tau}$. Thus, the distribution remains the same form under a change of scale.\\

Consequently, distributions of avalanches properties such as number of topplings, affected area, and radius of gyration may exhibit power-law behaviours over a truncated range, as widely studied by many research \cite{BTW1988, Lubeck1997Numerical,Yadav2022FiniteSize}. Notably, on a finite lattice, the power-law behaviour is truncated by a cutoff dependent on system size. The finite size scaling (FSS) hypothesis represents this systematically as 

\begin{equation}\label{eq:FSSH}
P(x,L)\sim x^{-\tau_x} f(\frac{x}{L^{D_x}})
\end{equation}
where $\tau_x$ is the critical exponent and the function $f$ is the finite-size scaling function that decides the cutoff location \cite{ahmed2010fss}.\\

Nevertheless, accurate estimation for these critical exponents remain elusive as research findings are not uniform. An early research by L\"{u}beck and Usadel \cite{Lubeck1997Numerical} examined the behaviour of probability distributions of topples ($s$), area ($a$) and and radius of gyration ($r$), and found estimates $\tau_s \approx 1.293$, $\tau_a \approx 1.330$ and $\tau_r \approx 1.665$ respectively. Meanwhile, recent work by Yadav et al. in 2022 \cite{Yadav2022FiniteSize} has further suggested that the area distribution contains a logarithmic dependence on system size. That is, they rescaled the probability $p(x)$ by $\log N$ where $N=L^2$ and argued the estimation should be closer to $\tau_a \approx 1$.\\

In addition, whether the BTW and Manna models belong to the same universality class remains a debated question. The Manna model replaces the deterministic rule in BTW with stochastic toppling \cite{Manna1991TwoState}, providing a natural extension to observe and compare their behaviours. Early work by Pietronero et al. \cite{pietronero1994} using normalisation group classified two models under the same universality class and characterised them by same exponents, with numerical estimates $\tau_a \approx 1.253$ and $\tau_a \approx 1.234$ for BTW and Manna respectively. Work by Chessa et al. \cite{Chessa1999Universality} with data collapse and moment analysis supported this classification and estimated $\tau_a \approx 1.27$ for both models. Nevertheless, other numerical studies have reached a different conclusion. L\"{u}beck and Usadel \cite{Lubeck1997Numerical} found differences in $\tau_a$ and $\tau_r$ between the BTW and Manna models that cannot be explained by error bars. Subsequent research by Biham et al. \cite{biham2001evidence} similarly argued deterministic Abelian models and stochastic models indicate different scaling behaviours and hence universality classes.\\

Therefore, motivated by these contrasting findings, we mainly explore two questions in this report: 
\begin{enumerate}
    \item What distributions do topples, area, loss and length exhibit, and for those that exhibit power-law behaviours with FSS, does the $\log N$ scaling proposed by Yadav et al. \cite{Yadav2022FiniteSize} improve data collapse?
    \item How do results differ between the BTW sandpile model and Manna model?\\
\end{enumerate}

The plan of this report is as follows. Chapter \ref{ch:background} discusses the fundamental properties of the BTW model, such as the role of open boundaries and the Abelian property. We also seek to highlight the difference between Manna's stochastic toppling rule and the deterministic toppling rule in BTW model. Chapter \ref{ch:methodology} explains simulations design and methods we use to measure the fits.  Chapter \ref{ch:results} presents the numerical results of our analysis. And lastly Chapter \ref{ch:discussion} interprets our findings and reflects on limitations.\\


